% Draft introduction and related work for downstream rewriting by writing-agent.
% Every citation key below appears in references.bib.

\section{Introduction}

Messenger RNA has moved rapidly from a niche reagent to a central modality
for vaccines, protein replacement, and \emph{in situ} gene editing. The
nucleoside-modified platform validated during the COVID-19 pandemic
\cite{thomas2021bnt162b2,karik{\'o}2013pseudouridine} and the arrival of
individualized cancer vaccines \cite{sahin2026tnbcvaccine} demonstrate that
synthetic mRNAs can direct therapeutically meaningful expression in humans.
Parallel advances in lipid nanoparticle formulation are extending these
gains to non-vaccine cargos, including proteins encoded at
low doses for tissue-specific applications
\cite{kim2026lnp_utr,wisitrasameewong2022mrnaregen}. Yet the \emph{sequence}
of the mRNA itself, and in particular its 5' untranslated region, remains
a surprisingly under-optimized lever.

The 5' UTR is a principal determinant of ribosome recruitment and
start-codon selection, and even small changes to it can shift the amount
of protein produced by orders of magnitude. Strong Kozak context, scanning
efficiency, upstream AUGs, and local secondary structure all contribute
to initiation \cite{chaldebas2026utrvariants,stettler2026kozak,
balasov2025kozaknoncanonical,shikalov2021ddx3}, but their combined
effect is difficult to predict from first principles. Massively parallel
reporter assays (MPRAs) coupled to polysome profiling have emerged as a
practical way to quantify 5' UTR activity at scale
\cite{reimaopinto2025danio5p,strayer2024naptrap}, yielding datasets
large enough to train supervised models that predict mean ribosome load
(MRL) from sequence alone. A longstanding limitation of these efforts is
that most reporter libraries have been profiled in a single cell type,
typically HEK293T, and it has not been established whether models trained
on such data generalize to therapeutically relevant cells such as primary
T lymphocytes or liver-derived lines.

Sequence design has advanced in parallel. Gradient-based activation
maximization over one-hot encodings, exemplified by Fast SeqProp
\cite{linder2021fastseqprop}, produces high-scoring sequences but is
prone to local optima and limited diversity. Generative approaches, in
particular Deep Exploration Networks (DENs)
\cite{linder2020den}, address both issues by training a generator to
maximize a fitness predictor while penalizing similarity between outputs;
variational autoencoders (VAEs) can additionally constrain generated
sequences to stay within the likelihood region of the training data.
These ideas have been applied to polyadenylation signals, splicing
regulators, and protein design
\cite{linder2020den,lei2025promoterdiffusion,zhao2024promotergan},
and recently to m$^1$$\Psi$-modified 5' UTRs for vaccine use
\cite{tang2023smart5utr,kim2026lnp_utr}. Whether DEN- and SeqProp-designed
5' UTRs carry functional benefit beyond reporter readouts, and whether
they improve delivery of clinically relevant cargos such as mRNA-encoded
gene editors, is still an open question.

Gene editing by mRNA-delivered nucleases is a particularly demanding
test case. Editing efficiency depends non-linearly on transient
protein concentration, and the payloads are large and structurally
idiosyncratic. megaTAL nucleases, hybrid enzymes built from LAGLIDADG
homing endonuclease scaffolds and TALE arrays
\cite{mcmurrough2018laglidadg,laforet2019endonuclease,gwiazda2016tcellediting},
offer high on-target activity and have been used to disrupt
therapeutically relevant loci such as \emph{TGFBR2} in
tumor-infiltrating lymphocytes \cite{fix2022tgfbr2} and \emph{PDCD1}
in adoptive T-cell products \cite{mohammad2025pd1}. Delivery of
megaTALs as mRNA, via electroporation
\cite{schwarze2021electroporation,gwiazda2016tcellediting}, lets the
nuclease act transiently but leaves editing rates highly sensitive to
translation efficiency in the target cell.

In this work, we close the loop between MPRA-based modelling and
mRNA-delivered gene editing. We profile randomized 5' UTR libraries in
three human cell types by polysome-profiling MPRA, train convolutional
and k-mer models on the resulting data, and use them to guide the
design of 50-nucleotide and 25-nucleotide 5' UTRs via Fast SeqProp and
DENs. We then test designed 5' UTRs in a mRNA-delivered megaTAL gene
editing assay targeting two independent loci in two cell lines. Our
central contributions are: (i) empirical demonstration that 5' UTR
activity is highly correlated across cell types and therefore transferable;
(ii) a new model, Optimus 5-Prime(25), adapted to 25-nt fully-variable
libraries; and (iii) functional validation that generatively-designed
5' UTRs support absolute editing efficiencies above 40\% for
\emph{TGFBR2} and above 80\% for \emph{PDCD1} at therapeutic mRNA doses,
with cargo- and locus-specific deviations that argue for empirical
validation alongside \emph{in silico} design.

\section{Related Work}

\paragraph{Deep-learning models of 5' UTR translation.} MPRA--trained
convolutional and attention models have mapped sequence features that
govern ribosome recruitment, including length, upstream ORFs, Kozak
context, and motif composition
\cite{reimaopinto2025danio5p,strayer2024naptrap,chaldebas2026utrvariants}.
Most such efforts profile a single cell type, leaving the cross-cell-line
transferability of predictions unresolved.

\paragraph{Generative sequence design.} Gradient-based activation
maximization (Fast SeqProp) and generative networks with diversity
penalties (DENs) \cite{linder2020den,linder2021fastseqprop} have become
the standard toolkit for cis-regulatory design. Related approaches
leveraging GANs and diffusion models have been applied to synthetic
promoters \cite{zhao2024promotergan,lei2025promoterdiffusion} and to
m$^1$$\Psi$-modified 5' UTRs for vaccines \cite{tang2023smart5utr}.

\paragraph{mRNA therapeutics.} Nucleoside-modified, LNP-formulated mRNA
underpins approved COVID-19 vaccines and clinical-stage individualized
cancer vaccines \cite{thomas2021bnt162b2,sahin2026tnbcvaccine,
karik{\'o}2013pseudouridine}. Recent work couples ionizable lipid
engineering with UTR tuning to lift potency at low doses
\cite{kim2026lnp_utr}.

\paragraph{megaTAL-based gene editing.} megaTALs combine engineered
LAGLIDADG homing endonucleases with TALE DNA-binding arrays and have
been deployed for high-efficiency editing of primary human T cells
\cite{gwiazda2016tcellediting,mcmurrough2018laglidadg,
laforet2019endonuclease}. TGFBR2 knockout lifts TGF-$\beta$-mediated
suppression of tumor-infiltrating lymphocytes \cite{fix2022tgfbr2} and
PDCD1 disruption is being pursued to relieve T-cell exhaustion
in CAR-T and TIL products \cite{mohammad2025pd1}. Electroporation of
synthetic mRNA is a standard delivery route for such edits
\cite{schwarze2021electroporation}, making the translational tuning of
the 5' UTR directly relevant to editing outcomes.
